% =========================================================
% Derivative Definition
% Chapter: Differentiation — Volume III
% =========================================================

\section{The Derivative}

% ---------------------------------------------------------
% Toolkit
% ---------------------------------------------------------
\begin{tcolorbox}[colback=gray!6, colframe=gray!40, arc=2pt,
  left=6pt, right=6pt, top=4pt, bottom=4pt,
  title={\small\textbf{Derivative Definition — Quick Reference}},
  fonttitle=\small\bfseries]
\small
\begin{tabular}{l l l l}
\toprule
\textbf{Label} & \textbf{Statement} & \textbf{Method} & \textbf{Detail} \\
\midrule
D\ref{def:derivative-at-point}
  & $f'(c) = \lim_{x\to c}\frac{f(x)-f(c)}{x-c}$
  & Definition
  & \hyperref[def:derivative-at-point]{$\downarrow$} \\
D\ref{def:differentiable-on-interval}
  & $f$ differentiable on $I$
  & Definition
  & \hyperref[def:differentiable-on-interval]{$\downarrow$} \\
T\ref{thm:differentiable-implies-continuous}
  & Differentiable $\Rightarrow$ continuous
  & Direct
  & \hyperref[thm:differentiable-implies-continuous]{$\downarrow$} \\
T\ref{thm:caratheodory-factorization}
  & Carathéodory factorization
  & Direct
  & \hyperref[thm:caratheodory-factorization]{$\downarrow$} \\
\bottomrule
\end{tabular}
\end{tcolorbox}

% ---------------------------------------------------------
\subsection{The Derivative at a Point}
% ---------------------------------------------------------

\begin{remark}[Motivation]
The derivative captures the instantaneous rate of change of $f$ at $c$.
The difference quotient
\[
\frac{f(x) - f(c)}{x - c}
\]
is the slope of the secant line through $(c, f(c))$ and $(x, f(x))$.
As $x \to c$, if this ratio converges to a finite limit, that limit is
the slope of the tangent line at $c$ --- the derivative $f'(c)$.

Two equivalent formulations appear in the literature.
The \emph{limit-at-$c$} form uses $x \to c$ as the primary variable.
The \emph{increment} form substitutes $h = x - c$ and lets $h \to 0$.
Both express the same limit; the increment form is more convenient
for algebraic manipulations, the limit-at-$c$ form for continuity arguments.
\end{remark}

\begin{tcolorbox}[colback=propbox, colframe=propborder, arc=2pt,
  left=6pt, right=6pt, top=4pt, bottom=4pt,
  title={\small\textbf{Definition (\texorpdfstring{%
    \hyperref[prf:derivative-at-point]{Derivative at a Point}}{Derivative at a Point})}},
  fonttitle=\small\bfseries]
\begin{definition}[\texorpdfstring{%
  \hyperref[prf:derivative-at-point]{Derivative at a Point}}{Derivative at a Point}]
\label{def:derivative-at-point}
Let $I \subseteq \mathbb{R}$ be an open interval and $f : I \to \mathbb{R}$.
Let $c \in I$.
The \emph{derivative of $f$ at $c$} is
\[
f'(c) \;:=\; \lim_{x \to c} \frac{f(x) - f(c)}{x - c},
\]
provided this limit exists as a finite real number.
Equivalently, setting $h = x - c$:
\[
f'(c) \;=\; \lim_{h \to 0} \frac{f(c + h) - f(c)}{h}.
\]
When $f'(c)$ exists, $f$ is said to be \emph{differentiable at $c$}.
\end{definition}
\end{tcolorbox}

\begin{remark*}[Fully quantified form]
$f$ is differentiable at $c$ with derivative $L$ if and only if there exists
$L \in \mathbb{R}$ such that for every $\varepsilon > 0$ there exists $\delta > 0$
such that for every $x \in I$ with $0 < |x - c| < \delta$,
\[
\left|\frac{f(x) - f(c)}{x - c} - L\right| < \varepsilon.
\]
\end{remark*}

\begin{remark*}[Predicate statement]
\[
\operatorname{DerivativeAt}(f,\,c,\,L)
\;:\iff\;
\operatorname{FunctionLimit}\!\left(\lambda x.\,\frac{f(x)-f(c)}{x-c},\;c,\;L\right).
\]
\end{remark*}

\begin{remark}[The two formulations are equivalent]
Setting $h = x - c$ converts one form to the other bijectively.
The condition $0 < |x - c| < \delta$ becomes $0 < |h| < \delta$;
$f$ being undefined at $h = 0$ is already excluded by $h \neq 0$.
\end{remark}

\begin{remark}[Standard notation]
$f'(c)$, $\;\dfrac{df}{dx}\big|_{x=c}$, $\;Df(c)$, and $\;\dot{f}(c)$
all denote the same quantity.
Leibniz notation $\dfrac{df}{dx}$ makes the chain rule visually transparent.
Prime notation $f'$ is most compact for real analysis.
\end{remark}

\begin{remark}[What the definition requires]
Three things are needed for $f'(c)$ to exist:
\begin{enumerate}
  \item $f$ is defined on an open interval containing $c$,
        so the difference quotient is defined for $x$ near $c$ on both sides.
  \item The left and right limits of the difference quotient as $x \to c$ agree.
  \item That common limit is finite.
\end{enumerate}
If the limit is $\pm\infty$, the function has a vertical tangent at $c$
and is not differentiable there in this sense.
\end{remark}

\begin{remark}[Canonical non-example: $|x|$ at $0$]
For $f(x) = |x|$:
\[
\frac{f(x) - f(0)}{x - 0} = \frac{|x|}{x} =
\begin{cases} +1 & x > 0, \\ -1 & x < 0. \end{cases}
\]
The one-sided limits disagree, so $f'(0)$ does not exist.
This is the canonical example that continuity does not imply differentiability.
\end{remark}

% ---------------------------------------------------------
\subsection{Differentiability on an Interval}
% ---------------------------------------------------------

\begin{tcolorbox}[colback=propbox, colframe=propborder, arc=2pt,
  left=6pt, right=6pt, top=4pt, bottom=4pt,
  title={\small\textbf{Definition (\texorpdfstring{%
    \hyperref[prf:differentiable-on-interval]{Differentiable on an Interval}}{Differentiable on an Interval})}},
  fonttitle=\small\bfseries]
\begin{definition}[\texorpdfstring{%
  \hyperref[prf:differentiable-on-interval]{Differentiable on an Interval}}{Differentiable on an Interval}]
\label{def:differentiable-on-interval}
Let $I \subseteq \mathbb{R}$ be an interval and $f : I \to \mathbb{R}$.
\begin{itemize}
  \item $f$ is \emph{differentiable on the open interval $(a,b)$} if $f$ is
        differentiable at every $c \in (a,b)$.
  \item $f$ is \emph{differentiable on the closed interval $[a,b]$} if $f$ is
        differentiable on $(a,b)$, the right derivative at $a$ exists, and the
        left derivative at $b$ exists.
\end{itemize}
The \emph{derivative function} $f' : I \to \mathbb{R}$ assigns to each
differentiable point its derivative value.
\end{definition}
\end{tcolorbox}

\begin{remark*}[Fully quantified form]
$f$ is differentiable on $(a,b)$ if and only if for every $c \in (a,b)$,
$\operatorname{DerivativeAt}(f,\,c,\,f'(c))$ holds.
\end{remark*}

\begin{remark*}[Predicate statement]
\[
\operatorname{DifferentiableAt}(f,\,c)
\;:\iff\;
\exists L \in \mathbb{R},\;\operatorname{DerivativeAt}(f,\,c,\,L).
\]
\end{remark*}

% ---------------------------------------------------------
\subsection{Differentiable Implies Continuous}
% ---------------------------------------------------------

\begin{theorem}[\texorpdfstring{%
  \hyperref[prf:differentiable-implies-continuous]{Differentiable Implies Continuous}}{Differentiable Implies Continuous}]
\label{thm:differentiable-implies-continuous}
Let $f : I \to \mathbb{R}$ and $c \in I$.
If $f$ is differentiable at $c$, then $f$ is continuous at $c$.
\end{theorem}

\begin{remark*}[Fully quantified form]
For every open interval $I$ and every $c \in I$, if $f$ is differentiable at $c$
then $\lim_{x \to c} f(x) = f(c)$.
\end{remark*}

\begin{remark*}[Predicate statement]
\[
\operatorname{DifferentiableAt}(f,\,c)
\;\Rightarrow\;
\operatorname{ContinuousAt}(f,\,c).
\]
\end{remark*}

\begin{remark}[Proof idea]
Write $f(x) - f(c) = \dfrac{f(x)-f(c)}{x-c} \cdot (x-c)$.
As $x \to c$: the difference quotient converges to $f'(c)$,
and $(x-c) \to 0$. Their product therefore converges to $f'(c) \cdot 0 = 0$,
giving $f(x) \to f(c)$.
\end{remark}

\begin{remark}[Converse fails]
Continuity does not imply differentiability.
$f(x) = |x|$ is continuous everywhere but fails to be differentiable at $0$.
More dramatically, the Weierstrass function is continuous everywhere
and differentiable nowhere.
\end{remark}

\begin{remark}[Position in the regularity hierarchy]
\[
\text{differentiable at } c
\;\Rightarrow\;
\text{continuous at } c
\;\not\Rightarrow\;
\text{differentiable at } c.
\]
Every assumption of differentiability is strictly stronger than an
assumption of continuity alone.
\end{remark}

% ---------------------------------------------------------
\subsection{Carathéodory Factorization}
% ---------------------------------------------------------

\begin{theorem}[\texorpdfstring{%
  \hyperref[prf:caratheodory-factorization]{Carathéodory Factorization}}{Caratheodory Factorization}]
\label{thm:caratheodory-factorization}
Let $f : I \to \mathbb{R}$ and $c \in I$.
Then $f$ is differentiable at $c$ with $f'(c) = L$ if and only if there
exists a function $\varphi : I \to \mathbb{R}$ such that
\begin{enumerate}
  \item $\varphi$ is continuous at $c$,
  \item $\varphi(c) = L$, and
  \item $f(x) - f(c) = \varphi(x)(x - c)$ for all $x \in I$.
\end{enumerate}
\end{theorem}

\begin{remark*}[Fully quantified form]
$f$ is differentiable at $c$ with derivative $L$ if and only if there exists
$\varphi : I \to \mathbb{R}$, continuous at $c$ with $\varphi(c) = L$,
such that $f(x) - f(c) = \varphi(x)(x-c)$ for all $x \in I$.
\end{remark*}

\begin{remark*}[Predicate statement]
\[
\operatorname{DerivativeAt}(f,\,c,\,L)
\;\iff\;
\exists\,\varphi : I \to \mathbb{R},\;
\bigl(
\operatorname{ContinuousAt}(\varphi,\,c)
\;\land\;
\varphi(c) = L
\;\land\;
\forall x \in I,\; f(x) - f(c) = \varphi(x)(x-c)
\bigr).
\]
\end{remark*}

\begin{remark}[The explicit factorizing function]
The unique $\varphi$ satisfying the identity is
\[
\varphi(x) =
\begin{cases}
\dfrac{f(x) - f(c)}{x - c} & x \neq c, \\[6pt]
L & x = c.
\end{cases}
\]
Differentiability of $f$ at $c$ is precisely the statement that this
piecewise-defined function is continuous at $c$.
The factorization eliminates the division-by-zero that makes the
limit definition awkward in algebraic arguments.
\end{remark}

\begin{remark}[Chain rule via Carathéodory]
If $f$ is differentiable at $c$ with factor $\varphi$, and $g$ is
differentiable at $f(c)$ with factor $\psi$, then
\begin{align*}
g(f(x)) - g(f(c))
  &= \psi(f(x))\cdot(f(x) - f(c)) \\
  &= \psi(f(x))\cdot\varphi(x)\cdot(x - c).
\end{align*}
The function $x \mapsto \psi(f(x))\cdot\varphi(x)$ is continuous at $c$
and evaluates there to $g'(f(c))\cdot f'(c)$.
This gives the chain rule with no limit computation.
\end{remark}

\begin{remark}[Why this matters]
The Carathéodory form is the preferred tool for proving the product rule,
quotient rule, and chain rule cleanly. It converts limit arguments into
continuity arguments, which are often simpler to chain together.
\end{remark}
